\section{系统辨识：从输入输出数据估计动力学}
\label{sec:16-Control-and-Reinforcement-Learning/01-Optimal-Control/07}

前面五节每一节都从 \(\dot x = f(x,u)\) 或 \(x_{t+1}=f(x_t,u_t)\) 起步，把它当成已知条件。可这个 \(f\) 从哪来？供暖系统的房间热容、墙体的蓄放热时间常数、室外温度的渗入系数，说明书上没有，供应商也说不清。真实做法是：装上传感器，让系统跑一段，从记录下来的输入输出里把 \(f\) 估出来。这一步一动手，控制问题就变成了统计问题，统计学的每一样毛病都跟着进来。

\subsection{已知状态时，辨识就是一次最小二乘}

从最简单的形式开始。线性离散模型

\[
x_{t+1} = A x_t + B u_t + w_t
\]

里待估的是 \((A,B)\)。若状态全部可测，把若干段数据按列排开，记 \(X_-\) 为各时刻的状态、\(U_-\) 为对应的输入、\(X_+\) 为下一时刻的状态：

\[
X_+ = \begin{pmatrix} A & B\end{pmatrix}
\begin{pmatrix} X_- \\ U_- \end{pmatrix} + W .
\]

这是一个多输出线性回归。解法是熟悉的：

\[
\begin{pmatrix}\widehat A & \widehat B\end{pmatrix}
= X_+ Z^{\trans}\bigl(Z Z^{\trans}\bigr)^{-1},
\qquad
Z = \begin{pmatrix} X_- \\ U_-\end{pmatrix}.
\]

计算这一步没有难度。难度全在那个 \((ZZ^{\trans})^{-1}\) 上——它能不能求逆、条件数好不好，取决于你的实验怎样激发系统，而不取决于你的算法。

\subsection{没有持续激励，参数在数学上就分不开}

设想供暖一直开在固定档位。房间温度在一个窄区间里小幅晃动，记录下来的 \(Z\) 各行几乎共线，\(ZZ^{\trans}\) 接近奇异。此时有一整族 \((A,B)\) 把观测数据解释得一样好，最小二乘挑出来的那个只是被噪声选中的一个。这是可辨识性的失效，不是样本量不够——再记十年数据也补不上。

条件写下来叫持续激励：输入的谱要在你关心的那些频率上有足够能量，使 \(ZZ^{\trans}\) 在相关方向上有界远离奇异。它与 \hyperref[sec:11-Mathematical-Statistics/03-Sufficient-Statistics-and-Information/02]{``Score 与 Fisher 信息衡量局部可辨认性''} 说的是同一件事的两个说法——信息矩阵在某个方向上退化，那个方向的参数就估不动。

激励也不是越猛越好。实体系统受舒适度、安全余量、执行器磨损与能耗的约束，你不能为了辨识把机房温度拉到四十度。折中的手法有几种：小幅多频的伪随机信号，把能量摊薄在宽频带上；分阶段实验，每段只激发一部分模态；或者退而使用运行中的自然波动。哪一种都要在报告里写清楚，哪些动态方向从头到尾没被激发过。参数的置信区间很窄，救不了一个在结构上就辨认不出来的方向。

\subsection{闭环数据把控制动作和噪声绑在一起}

真实系统很少能停下反馈单独做实验，数据往往是控制器在线跑着的时候顺手记的。这时 \(u_t\) 由过去的读数决定，而读数里含着噪声与扰动。于是回归自变量与误差项相关，最小二乘的无偏性前提当场破掉。

后果具体而有方向性。控制器为抵消扰动而做的反应，会被回归当成被控对象的动力学记到 \(A\)、\(B\) 头上；辨识出来的模型于是带着控制器的影子，换一个控制器就不成立了。这与观察数据里的混杂是同一个结构——动作依据状态选择，状态又受扰动影响，扰动同时进入下一时刻的状态，形成一条后门路径，\hyperref[sec:13-Machine-Learning/12-Causal-Inference/03]{``混杂、碰撞点与后门调整''} 里画的图搬过来就能用。

应对手段各自处理这条依赖。工具变量找一个与噪声无关、只通过 \(u\) 影响输出的外生激励，比如叠加在控制信号上的独立抖动；预测误差法直接把控制器写进似然，估计条件于整个反馈结构；再不然就在实验期间断开反馈，用开环数据辨识，代价是这段时间系统不受保护。

\subsection{只有输出时，状态坐标本身也是未定的}

更常见的情形是状态测不全，能看见的只有

\[
y_t = C x_t + v_t .
\]

此时 \(x_t\) 与 \((A,B,C)\) 都是未知量。子空间辨识用块 Hankel 矩阵组织数据，靠低秩结构一次性恢复出一个等价的状态表示；期望最大化则交替进行，E 步用 Kalman 平滑给出隐状态的后验，M 步据此更新参数，用到的正是 \hyperref[sec:13-Machine-Learning/07-Graphical-and-Sequential-Models/03]{``Kalman filter 是动态系统上的 Gaussian 条件化''} 那套递推。

这里有一条结构性的不确定，跟数据量无关。取任意可逆矩阵 \(S\)，令 \(z_t = S^{-1}x_t\)，则

\[
(A, B, C) \longmapsto (S^{-1}AS,\; S^{-1}B,\; CS)
\]

给出完全相同的输入输出行为。状态坐标只确定到相似变换。所以拿到 \(\widehat A\) 的某一行去说``这一维代表墙体温度''是没有根据的，除非你另外施加了固定坐标的约束。可以放心解读的是相似变换下的不变量：特征值决定的时间尺度、稳定模态、静态增益、可控可观的秩。

\subsection{辨识误差进入控制回路会被放大}

到这里为止都还是统计。把系统辨识与一般回归区分开的，是估出来的模型接下来要被拿去设计控制器。

先看放大的机制。前面 LQR 给出的最优增益依赖 \(A, B\)，闭环矩阵是 \(A - BK\)。\(A\) 的估计偏了 \(\Delta\)，实际闭环变成 \(A + \Delta - BK\)；\(K\) 是照着 \(\widehat A\) 算出来的，它把闭环极点尽可能推向左边，推得越狠，极点对 \(\Delta\) 的敏感度越高。为了追求名义代价最小而把增益调激进，等于把模型误差的杠杆一起调大。名义性能与鲁棒性在这里是对着干的两件事。

再看数据分布的问题。辨识数据是在某个旧控制器（或开环实验）下采出来的，覆盖了状态空间的某一片。新控制器上线后会把系统带到别处——它正是为了做得更好才改变轨迹的。模型在没有数据支撑的区域上外推，误差没有任何东西约束。

\begin{center}
\begin{tikzpicture}[>=stealth, every node/.style={font=\small}]
  \draw[->,gray!70] (-0.3,0) -- (8.2,0) node[right,black] {\(x_1\)};
  \draw[->,gray!70] (0,-0.3) -- (0,3.7) node[above,black] {\(x_2\)};

  \draw[dashed,gray!70] (2.2,1.3) ellipse (1.9 and 0.8);
  \foreach \p in {(1.1,1.2),(1.5,1.55),(1.6,0.95),(2.0,1.4),(2.2,1.05),(2.4,1.65),
                  (2.7,1.25),(3.0,1.0),(3.1,1.55),(3.4,1.3),(1.9,1.15),(2.6,1.5)}
    \fill[gray!65] \p circle (1.6pt);
  \node[gray!55!black,anchor=north] at (2.2,0.35) {辨识数据覆盖的区域};

  \draw[very thick,blue!55!black] plot[smooth] coordinates
    {(1.2,1.25) (2.1,1.55) (3.1,1.95) (4.1,2.35)};
  \draw[very thick,red!70!black,dashed] plot[smooth] coordinates
    {(4.1,2.35) (5.2,2.75) (6.3,2.85) (7.2,2.55)};
  \fill[blue!55!black] (1.2,1.25) circle (2pt);
  \fill[red!70!black] (7.45,2.5) circle (2.2pt);
  \node[anchor=west] at (7.55,2.5) {目标};
  \node[red!70!black,anchor=south] at (5.6,3.05) {模型在此处无依据};
  \node[blue!55!black,anchor=south] at (2.6,2.45) {新控制器的闭环轨迹};
\end{tikzpicture}

\emph{控制器为了改善性能，会把系统带出采集数据的那片区域；模型误差在图中虚线段上不再受任何数据约束。}
\end{center}

图里那段虚线解释了为什么``预测精度高''不等于``可以拿来做控制''。一个在历史数据上残差极小的黑箱，可能只是把旧控制器的行为背了下来。

因此验收要按控制的用途来做，而不是按拟合优度。至少三件事该查。做自由仿真：不用一步预测，而是让模型从初始状态自己往前跑几十步，与真实轨迹比——一步预测的残差小得可以是因为下一时刻的状态本来就和当前差不多。查新息的白性：本单元讲 LQG 时用过的那把尺子在这里同样适用，残差若有自相关，说明模型阶数或噪声结构设错了，属于模型误设，而 \hyperref[sec:11-Mathematical-Statistics/06-Resampling-and-Model-Selection/04]{``偏差---方差、正则化与模型选择''} 里关于误设的讨论逐条适用。查参数的稳定性：换数据窗口、换模型阶数，估计值若大幅跳动，说明你落在了辨识数据支撑不住的复杂度上。

设计侧还要留出余量。把辨识结果当成一个模型集合而不是一个点估计——参数的置信椭球、或者一族有界摄动，然后要求控制器对整个集合都稳定。牺牲一点名义性能，换回图中虚线段上不至于失控。这条思路也解释了一个常见的经验：保留了关键慢模态和静态增益的简化模型，做起控制来常常比拟合得更好的复杂黑箱更安全。

至此这个单元的链条闭合了：有了模型才能算最优，模型来自数据，数据里的动作又来自控制器。链条上每一环都可能松动，而松动的方式已经逐条摆开。剩下的问题是当模型压根建不起来的时候怎么办——状态空间大到无法逐点求解，动力学复杂到无法写成方程，只能靠与环境反复交互来积累经验。\hyperref[ch:136]{``强化学习：模型未知时怎样逼近最优决策''} 从这里接手。
